\scanpage{68}
Comme, \uline{donnée}, c'est un hom de bifoncteurs
\[
F(X)\otimes\deleted{\unsure{*}}F(Y)\longrightarrow F(X\deleted{\otimes}*Y)
\]
or il suffit de \deleted{\unsure{le}} le donner sur la sous catégorie \inserted{\unsure{de la|$k$-lin.}} $\Delta\times\Delta$ \inserted{de $\mathbb{K}\times\mathbb{K}$} des $\mathbb{C}^i\times \mathbb{C}^j$, (avec sa structure \unsure{$k$-lin|bi-lin}) donc de se donner des hom de bicomplexe de cochaînes
\[
\underset{C^{\bullet}}{\underset{\Vert}{F(\mathbb{C}^{\bullet})}}\otimes_B
\underset{C^{\bullet}}{\underset{\Vert}{F(\mathbb{C}^{\bullet})}}\longrightarrow
F(\deleted{\unsure{\Phi}}\mathbb{C}^{\bullet}*\deleted{\unsure{P}}\mathbb{C}^{\bullet})\quad\text{dans } B^{\bullet\bullet}
\]
ce qu'on \unsure{peut encore|pourrait|peut} expliciter, si $\mathbb{C}^i\deleted{\otimes}*\mathbb{C}^j$ correspond au bicomplexe de chaînes \inserted{\unsure{flasques de|défini par}} $\mathbb{C}^i_{\bullet}*\mathbb{C}^j_{\bullet}=\deleted{\unsure{C^i_{\bullet}\otimes}}\operatorname{ND}_{\bullet}(\underset{\bigwedge^{i+1}\Phi_*\otimes\bigwedge^{j+1}\Phi_*}{\underset{\Vert}{\mathbb{C}^i_*\otimes \mathbb{C}^j_*}})$
% slanted
\margin{cas le plus commun: $B=\unsure{\mathit{Ab}_k|\mathit{Ab}}$}
\[
\underset{C^i}{\underset{\Vert}{F(\mathbb{C}^i)}}\otimes_B
\underset{C^j}{\underset{\Vert}{F(\mathbb{C}^j)}}\longrightarrow
\operatorname{Hom}^{\bullet}_k\bigl(\underset{C_i^{\bullet}\deleted{\otimes}*C_j^{\bullet}}{\underset{\Vert}{(\mathbb{C}^i_{\bullet}*\mathbb{C}^j_{\bullet})^{\vee}}},C^{\bullet}\bigr)
\tag*{1)}
\]
i.e.
\[
(C^i_{\unsure{\xi|\zeta}}\otimes_B C^j)\otimes_k
\underbrace{\overbrace{(C_i^{\bullet}*C_j^{\bullet})}^{\mathrm{déf}\ \Vert\ C^{\bullet}_{i,j}\deleted{\unsure{N}}}}_{\in\mathbb{K}^{\bullet}}\longrightarrow C^{\bullet}
\tag*{2)}
\]
avec compatibilité aux op. diff. pour $i,j$ variables.

On peut aussi l'écrire
\[
\deleted{(C^i\otimes_B C^j)\otimes_k}
\]
\[
C^i\otimes_B C^j\longrightarrow
\operatorname{Hom}^{*}_{k}(\underbrace{C_i^*\otimes C_j^*}_{\bigwedge^{i+1}\Phi^*\otimes\bigwedge^{j+1}\Phi^*},\underset{\operatorname{DP}^*(C^{\bullet})}{\underset{\Vert}{C^*}})
\tag*{(1*)}
\]
ou encore
\[
(C^i\otimes_B C^j)\otimes_k
\underbrace{\bigl(\bigwedge^{i+1}\underline{\Phi}^{*}\overset{*}{\otimes}\bigwedge^{j+1}\underline{\Phi}^{*}\bigr)}_{C_{i,j}^{*}\unsure{\simeq}\,}
\xrightarrow{u^{i,j}}C^{*}
\tag*{(2*)}
\]
% archive sheet number: 67
